% !TEX program = lualatex
\documentclass[12pt, a4paper]{article}

\usepackage{fontspec}
\setmainfont{TeX Gyre Pagella}

\usepackage{amsmath, amssymb, amsthm}
\usepackage{geometry}
\usepackage{microtype}
\usepackage{setspace}
\usepackage{titlesec}
\usepackage{hyperref}
\usepackage{mdframed}

\geometry{
  top=1.25in,
  bottom=1.25in,
  left=1.35in,
  right=1.35in
}

\onehalfspacing

\hypersetup{
  colorlinks=true,
  linkcolor=black,
  citecolor=black,
  urlcolor=black
}

\newtheoremstyle{named}
  {12pt}{12pt}{\itshape}{}{\bfseries}{.}{.5em}{}
\theoremstyle{named}
\newtheorem{proposition}{Proposition}
\newtheorem{definition}{Definition}

\newmdenv[
  linewidth=0.8pt,
  innertopmargin=10pt,
  innerbottommargin=10pt,
  innerleftmargin=14pt,
  innerrightmargin=14pt,
  skipabove=14pt,
  skipbelow=14pt
]{keybox}

\titleformat{\section}{\large\bfseries}{\thesection.}{0.5em}{}
\titleformat{\subsection}{\normalsize\bfseries}{\thesubsection.}{0.5em}{}

\title{%
  \textbf{The Armillary Sphere and the Thought Experiment}\\[0.4em]
  \large On Proposals That Are Neither Blueprints Nor Metaphors
}
\author{Flyxion}
\date{}

% ---------------------------------------------------------------
\begin{document}

\maketitle
\thispagestyle{empty}

\begin{abstract}
Speculative proposals — imagined cities, hypothetical organisms,
impossible machines, impossible institutions — are routinely
misclassified. They are read either as blueprints (literal plans for
construction) or as metaphors (ornamental illustrations of something
already known). This essay argues that a significant class of
speculative proposals belongs to neither category. Following the logic
of the thought experiment in physics and the armillary sphere in
navigation, these proposals function as coordinate transformations:
instruments for making visible the assumptions that ordinarily remain
hidden inside our frameworks for understanding civilization, design,
ecology, and infrastructure. The essay develops a formal apparatus for
this claim, representing civilizations as points in a five-dimensional
parameter manifold $\mathcal{C}$ and armillary proposals as directional
vectors in the tangent space $T_C\mathcal{C}$. A visibility functional
$V(T)$ measures the epistemic value of a proposal by the significance
of the assumptions it makes visible rather than the feasibility of the
configuration it imagines. The Invariant Discovery Principle identifies
the true objects of study: not optimal civilizational configurations but
structural invariants that persist under broad families of
transformations. The Armillary Principle closes the formal account:
the value of a proposal is determined by $\Delta I$ and $\Delta\Omega$
— new invariants revealed and new reachable possibilities exposed —
not by the magnitude of its departure from present reality. The essay
concludes by identifying five civilizational axes — maintenance,
ecology, persistence, accessibility, and coordination — as the
invariant structure that a large body of otherwise disparate speculative
work collectively orbits.
\end{abstract}

\newpage
\tableofcontents
\newpage

% ---------------------------------------------------------------
\section{The Satirist's Error}

There is a reliable genre of online commentary that applies to
speculative proposals. The format is consistent. A list of ambitious,
unusual, or technically extreme ideas is assembled and presented in the
flattest possible prose — the surface features rendered faithfully, the
structural purpose systematically ignored. The result reads as satire
even when it is not intended to be. The ideas sound absurd because
abstraction has been collapsed into literalism: the thought experiment
has been read as a building permit.

The genre is not without value. It exposes a real risk in speculative
thinking — the risk of mistaking a conceptual probe for an engineering
proposal, of treating a coordinate transformation as a destination. But
it also systematically misreads what it is satirizing, because it
applies a single interpretive lens — \emph{is this a plan?} — to
objects that require a different question entirely.

Consider what twentieth-century civilization would look like if
described in the same register:

\begin{keybox}
The plan is to cover the planet with black rivers of asphalt, fill them
with exploding metal boxes powered by liquefied remains of ancient
organisms, connect everything with copper nerves suspended from wooden
poles, then teach billions of people to stare into glowing glass
rectangles for the majority of their waking hours.
\end{keybox}

This description is accurate. It sounds ridiculous. The ridiculousness
is not a property of the civilization but of the descriptive register
applied to it. Existing systems are not usually described at the level
of their physical substrates without their functions and purposes
supplied. Speculative systems almost always are. The asymmetry is the
source of most satirical energy directed at ambitious proposals.

The deeper error is epistemic. Treating all proposals as candidates
for literal implementation misses the category of proposals whose
primary function is not to be built but to be thought. The armillary
sphere is the instrument for understanding this category.

% ---------------------------------------------------------------
\section{The Armillary Sphere}

An armillary sphere is a model of the heavens consisting of a nest of
rings representing celestial circles — the equator, the ecliptic, the
meridian, the tropics — arranged around a small central Earth. It was
used for teaching, for computation, and for navigation across more than
fifteen centuries of astronomical practice. It was carried on ships
and installed in observatories and given as gifts to rulers.

The armillary sphere is not a model of reality. The celestial circles
it represents do not exist as physical objects. The equator is not a
ring hanging in space; the ecliptic is not a track the Sun runs along.
These are coordinate systems — constructs that organize observation and
make certain calculations possible. The armillary sphere makes those
coordinate systems tangible. It is a model of a model, an instrument
for operating with abstractions that would otherwise remain merely
abstract.

What the armillary sphere provides is not a picture of the world but a
framework within which the world can be navigated. Its rings do not
represent features of reality; they represent the relationships between
features of reality — the angles, the cycles, the periodicities that
allow a navigator to situate a ship in the ocean or an astronomer to
predict a solstice. The instrument is useful not because the rings are
real but because the relationships they encode are real.

A civilization-scale thought experiment functions the same way. The
proposal for a city built from living forest structures is not primarily
a claim that such a city should or could be built. It is a ring on an
armillary sphere. It encodes a relationship: the relationship between
construction and ecology, between the timescales of human habitation and
the timescales of biological growth, between infrastructure as a fixed
artifact and infrastructure as an ongoing process. Rotating this ring
relative to others makes certain things visible that were not visible
before.

% ---------------------------------------------------------------
\section{Proposals as Coordinate Transformations}

In mathematics, a change of coordinates does not alter the underlying
object. It alters which features of the object are easy to see and which
are hidden. A circle expressed in Cartesian coordinates is the set of
points satisfying $x^2 + y^2 = r^2$; in polar coordinates, it is the
trivially simple $r = \text{const}$. The coordinate transformation does
not change the circle. It changes what is visible about the circle.

Speculative proposals function similarly. They change coordinates.

\begin{definition}[Coordinate Transformation Proposal]
A coordinate transformation proposal is a speculative design,
institutional arrangement, or technological system whose primary
epistemic function is to vary parameters that are ordinarily held fixed
in conventional analysis, thereby making visible the assumptions
embedded in conventional arrangements.
\end{definition}

When someone proposes habitation modules that float on the ocean, they
are varying a parameter — the fixity of land as the substrate of
civilization — that is so rarely varied that it has become invisible as
a parameter. It appears to be a natural feature of civilization rather
than a contingent choice. The proposal reveals that it is a choice by
imagining what follows if it is made differently.

When someone proposes that cities be organized around biological nutrient
cycles rather than economic sectors, they are varying the parameter that
treats ecology as an externality of civilization rather than its
substrate. The proposal does not primarily advance a plan; it asks a
question: what becomes visible if ecology is the primitive unit rather
than property?

When someone proposes that maintenance be the primary function of
infrastructure rather than construction, they are varying the parameter
that treats built environments as finished artifacts. Most infrastructure
planning proceeds as if the relevant question is \emph{what to build};
the proposal rotates this to \emph{what to sustain}, exposing the
degree to which maintenance is systematically undervalued by the
incentive structures of contemporary economies.

The individual proposals — yogurt-culturing robots exploring self-repair
mechanisms, SCOBY-based structural materials exploring low-energy
manufacturing, forest-integrated cities exploring ecological embedding,
modular floating habitats exploring flexible occupation, kelp cultivation
systems exploring carbon and nutrient cycling, distributed maintenance
organisms exploring repair infrastructure — are each rings on the
armillary sphere.

Each ring is a coordinate. Each explores a different axis along which
the assumptions of contemporary civilization can be varied. Each makes
visible a different set of contingencies that habituation has rendered
invisible.

Taken together they probe a single underlying question, which is neither
the question of any individual proposal nor their sum but the invariant
that all of them orbit:

\begin{keybox}
What does a civilization look like when maintenance, repair, persistence,
accessibility, ecology, and coordination are treated as first-class
design variables rather than externalities?
\end{keybox}

The individual proposals are not answers to this question. They are
instruments for approaching it.

\subsection{Civilizational Parameter Space}

To make the coordinate transformation language mathematically precise,
represent a civilization as a point in a five-dimensional configuration
space:
\[
  C = (m,\, e,\, p,\, a,\, c),
\]
where the coordinates are normalized in $[0,1]$ and denote:
\[
  m = \text{maintenance intensity}, \quad
  e = \text{ecological integration},
\]
\[
  p = \text{persistence orientation}, \quad
  a = \text{accessibility orientation}, \quad
  c = \text{coordination explicitness}.
\]
The set of possible civilizational configurations forms a manifold:
\[
  \mathcal{C} = \bigl\{(m,e,p,a,c) \in [0,1]^5\bigr\}.
\]
A coordinate transformation proposal acts as an operator
$T: \mathcal{C} \to \mathcal{C}$ which perturbs one or more
coordinates while holding others approximately fixed. Contemporary
industrial civilization occupies a region of $\mathcal{C}$ with low
$m$, low $e$, low $p$, low $a$, and low $c$: construction-dominant,
ecologically extractive, production-oriented, location-bound, and
coordination-emergent.

\begin{definition}[Armillary Proposal]
An armillary proposal specifies a direction vector
\[
  v \in T_C\mathcal{C}
\]
in the tangent space of civilizational configurations at the current
point $C$. The proposal's primary purpose is to explore the structural
consequences of moving in the direction $v$ — to reveal what features
of the current configuration are contingent rather than necessary.
An armillary proposal is not an optimal destination $C^\star$ but a
directional derivative: a probe into what the configuration manifold
looks like in a direction that conventional analysis does not explore.
\end{definition}

This formalizes the intuition that the yogurt robot is not about
yogurt, the kelp city is not about kelp, and the floating habitat is
not about floating. Each is a vector $v$ pointing away from the current
civilizational configuration along one of the five axes. The object
imagined is the arrow, not the destination.



Thought experiments in physics are the closest analog to what
speculative proposals at the civilizational scale are attempting.
Understanding how thought experiments work clarifies what speculative
proposals are doing and why the blueprint/metaphor dichotomy fails to
capture them.

No frictionless plane exists. No perfectly rigid body exists. No
adiabatic container is perfectly insulated. No Maxwell's demon hides
in any thermodynamic system. Schrödinger's cat is not an actual
laboratory protocol. The twin paradox does not describe two specific
people who have traveled in rockets. The thought experiments of physics
are populated by idealized objects that cannot be constructed and in
some cases cannot even be coherently imagined in full physical detail.

Yet these imaginary objects have been among the most productive epistemic
instruments in the history of science. Maxwell's demon forced a
reconceptualization of the relationship between information and
thermodynamics that eventually produced the modern understanding of
entropy, information theory, and computation. The twin paradox forced a
precise articulation of what simultaneity means in special relativity.
Schrödinger's cat forced a confrontation with the measurement problem
that is still generating productive work a century later. The
frictionless plane does not need to exist for Newton's laws to be stated
precisely; the idealization strips away complicating factors and reveals
the underlying structure.

What thought experiments share is the operation of selective idealization:
certain parameters are set to extreme or impossible values — friction to
zero, walls to perfect insulators, observers to perfect demons — in order
to see what the theory implies without the confounding effects of
realistic complexity. The result is not a description of a possible
experiment but a description of a possible world in which certain
features of the actual world are made maximally clear.

Civilizational thought experiments operate similarly but at a different
scale. Instead of setting friction to zero, they set ecology to primary.
Instead of making walls perfectly insulating, they make maintenance the
central function. Instead of introducing a perfect demon, they
introduce organisms whose purpose is repair rather than production.

\begin{proposition}[Thought Experiment Invariance]
The epistemic value of a thought experiment is not determined by the
feasibility of the imagined scenario but by whether the scenario
isolates a structural relationship that holds in the actual world and
is difficult to see without the idealization.
\end{proposition}

This can be made precise using a visibility functional. Let
$\mathcal{A}(C)$ denote the set of assumptions embedded in a
civilizational configuration $C$ — the parameters that are held fixed
and therefore invisible as parameters. For a transformation $T$, define:

\begin{definition}[Visibility Functional]
\[
  V(T) = \mu\!\bigl(\mathcal{A}(C) \setminus \mathcal{A}(T(C))\bigr),
\]
where $\mu$ measures the conceptual significance of the assumptions
made visible. An assumption becomes visible under $T$ when it was
treated as a constant in $C$ but is forced to vary in $T(C)$.
\end{definition}

\begin{proposition}[Visibility Principle]
The epistemic value of a coordinate transformation proposal is
proportional to $V(T)$ — the measure of previously invisible
assumptions it makes visible — rather than to the feasibility of
$T(C)$.
\end{proposition}

A proposal involving forty billion floating habitat modules is not
valuable because forty billion floating modules are a realistic
near-term objective. It is valuable if it isolates a structural
relationship — between occupation density, ecological throughput, energy
source, and maintenance infrastructure — that is difficult to see as
long as land-based urban settlement is treated as a fixed parameter.

The question to ask of any speculative proposal is not ``is this
feasible?'' but ``what does this make visible?''

\subsection{Model Organisms and Toy Problems}

Every mature field of inquiry maintains a collection of objects that
are valuable precisely because they are not the thing being studied.
These are intentionally simplified, idealized, or distorted objects
that expose hidden structure by removing the complexity that ordinarily
conceals it.

Biology uses \textit{Drosophila melanogaster} — the fruit fly — not
because geneticists want to understand fruit flies but because the
fly's short generation time, large chromosomes, and small genome make
visible relationships between genotype and phenotype that are present
in all organisms but difficult to observe in more complex ones. The
fly is a model organism: an intentionally chosen proxy that reveals
structural relationships at the cost of direct representational
accuracy.

Economists use the representative agent: a single maximizing individual
whose choices are taken to characterize the behavior of an entire
economy. The representative agent does not exist and no economist
believes it does. It is a model organism for economic analysis — a
simplification that makes certain equilibrium properties visible while
deliberately setting aside distributional complexity.

Computer scientists use toy problems: the traveling salesman, the
halting problem, the eight queens puzzle. These are not the problems
practitioners actually solve. They are chosen because their simplified
structure exposes algorithmic properties — complexity classes,
decidability boundaries, search space geometry — that are present in
real problems but obscured by irrelevant detail.

\begin{keybox}
A kelp city is to civilization what a frictionless plane is to
mechanics: an intentionally distorted object used to expose hidden
structure. The distortion is the point, not the defect.
\end{keybox}

Armillary proposals function as model organisms for civilizational
analysis. The yogurt robot is not proposed because self-repairing
yogurt cultures are the optimal maintenance technology. It is a model
organism for the relationship between biological processes and
mechanical infrastructure — a simplified object in which that
relationship is made visible without the confounding complexity of
actual maintenance systems. The floating city is a model organism for
the relationship between habitation, ecology, and land tenure — a
simplified object in which the contingency of land-based settlement
becomes observable.

The appropriate response to a model organism is not to evaluate whether
it is a good representation of the actual system. It is to ask what
structural properties it reveals. The fruit fly is a poor model of
mammalian development and an excellent model of genetic inheritance.
The kelp city is a poor model of near-term urban planning and an
excellent model of the relationship between occupation, nutrient cycling,
and carbon metabolism.



There is a genuine and serious criticism of ambitious speculative
proposals that survives the reframing offered above. It is the strongest
objection in the satirical genre, and it deserves engagement rather than
dismissal.

The objection is this: utopian technical proposals routinely assume that
replacing infrastructure automatically solves social problems. This
assumption is historically false.

A city made of living forest can still be oppressive. A city powered by
fusion can still be unequal. A city managed by artificial intelligence
can still be alienating. A city floating on the ocean can still become a
platform for extraction and enclosure. The technical substrate does not
automatically generate the social arrangements. History provides no
shortage of evidence. Garden cities were designed to solve industrial
alienation and became suburban dormitories. Brutalist housing was
designed to provide dignified space for working-class lives and became a
synonym for social failure. The technological utopianism of the early
internet promised a radical democratization of information and produced
surveillance capitalism and epistemic fragmentation.

The technical system and the social system are related but not identical.
Changing one does not automatically change the other.

This criticism lands, and it would land against any proposal that
implicitly treats technical redesign as social reform. But the armillary
sphere framing suggests a response. The proposals under consideration
are not primarily claiming that ecological infrastructure generates just
social arrangements. They are claiming something more specific and more
modest: that certain technical assumptions constrain the social
possibilities that remain reachable.

A civilization organized around extractive land use and disposable
infrastructure makes certain social arrangements very easy and others
very difficult. The goal of coordinate transformation proposals is not
to build a new utopia but to expand the set of reachable configurations
by varying the assumptions that currently constrain it. This is the
language of admissibility and reachability: the technical system does
not determine the social system, but it shapes which social systems
remain accessible.

\begin{keybox}
The question is not: does ecological infrastructure produce justice?

The question is: does disposable infrastructure foreclose it?
\end{keybox}

If certain technical arrangements make it structurally difficult for
communities to persist, to repair themselves, to maintain their
accumulated knowledge and social infrastructure across generations, then
the technical question is also a social question — not because the
technical determines the social but because the technical constrains the
reachable social space.

This is exactly the reachability analysis developed in institutional
contexts: institutions do not merely distribute resources within a fixed
space; they shape the space itself. The same logic applies to technical
infrastructure. The proposal for maintenance-first civilization is not
a plan for utopia; it is a proposal for expanding the admissible set of
social futures by treating continuity, repair, and ecological integration
as primary rather than derivative.

% ---------------------------------------------------------------
\section{The Intermediate Space}

The satirist cannot quite decide whether speculative proposals of this
kind are blueprints, jokes, manifestos, or philosophical thought
experiments. The indecision is correct. The proposals occupy a genuine
intermediate space that does not reduce to any of these categories.

They are not blueprints: the level of technical specification is
insufficient, and the purpose is not primarily construction. They are
not jokes: the underlying structural questions are serious and the
analysis motivating them is careful. They are not manifestos: they do
not primarily prescribe political action or assert moral imperatives.
They are not exactly thought experiments in the physics sense: they
involve multiple systems interacting in realistic complexity rather than
idealized single-parameter variation.

What they are is closer to what designers call provocations and what
philosophers call hypothetical imperatives: \emph{if} certain parameters
are varied, \emph{then} certain features of our current arrangements
become visible as contingent. The if is not a plan. The then is not a
promise. Together they constitute an invitation to think differently
about what is fixed and what is variable.

The armillary sphere captures this intermediate status precisely. An
armillary sphere is not a model of the Earth. It is not a toy. It is
not a work of art (though it can be beautiful). It is an instrument
for navigating a coordinate system that cannot be directly perceived.
The rings are real objects that represent abstract relationships. The
instrument is useful not when its rings are mistaken for real celestial
circles but when they enable a user to reason about real celestial
circles without being able to see them directly.

Civilizational proposals function the same way. The yogurt robot is a
real object in the imagination that represents an abstract relationship:
the relationship between biological processes and mechanical
maintenance, between living systems and manufactured ones, between
repair and production. The kelp farm is a real imagined object that
represents the relationship between food systems, carbon cycles, coastal
ecology, and distributed infrastructure. Each imagined object is a ring
on the sphere. The sphere is not the civilization; it is an instrument
for thinking about civilization.

% ---------------------------------------------------------------
\section{Ideas as Candidates for Exploration}

Most people treat ideas as candidates for belief. An idea is encountered,
evaluated against existing commitments, and either accepted or rejected.
The question asked of an idea in this mode is: is this true?

A different relationship to ideas treats them as candidates for
exploration. An idea is encountered, inhabited temporarily, followed
out into its implications, and then compared to other inhabited
frameworks. The question asked of an idea in this mode is not
\emph{is this true?} but \emph{what does this make visible?}

These are genuinely different activities. The first is appropriate for
claims about matters of fact — whether a particular bridge will hold,
whether a particular medication is effective, whether a particular
historical event occurred. The second is appropriate for frameworks —
coordinate systems, interpretive schemes, design principles, theoretical
vocabularies.

The confusion between these two activities generates most of the
misreadings of speculative work. When a proposal for floating cities is
evaluated as if it were a claim about matters of fact, the obvious
response is to ask whether floating cities are feasible, whether they
would be stable, whether people would want to live in them, whether the
economics work. These are reasonable questions about a blueprint. They
are the wrong questions to ask of a coordinate transformation.

The appropriate questions for a coordinate transformation proposal are:

What assumptions are being varied? What becomes visible when those
assumptions are released? What features of the current arrangement
appear contingent that appeared necessary before? What constraints
on social possibility are revealed as technical rather than natural?
What structural relationships does the imagined scenario isolate?

These questions cannot be answered by asking whether the proposal is
feasible. Frictionless planes are not feasible. They are among the most
productive thought experiments in the history of physics.

% ---------------------------------------------------------------
\section{The Invariant Structure}

If the individual proposals are rings on the armillary sphere, what is
the sphere itself? What is the invariant structure that all of them are
instruments for navigating?

The proposals listed above — self-repairing organisms, biological
construction materials, ecologically integrated settlements, modular
occupation units, nutrient-cycling aquaculture, distributed maintenance
infrastructure, wave-based memory architectures, admissibility-based
field theories — appear unrelated when examined at the level of their
surface content. Yogurt and field theory share no obvious domain.
Kelp and memory architectures share no obvious mechanism.

But at the level of the questions they probe, a common structure
emerges. Each proposal varies a parameter that contemporary industrial
civilization treats as fixed, and each varied parameter belongs to a
cluster:

\medskip

\noindent\textbf{Maintenance vs.\ construction.} Contemporary
infrastructure is organized around construction as the primary act and
maintenance as a secondary cost. The proposals collectively vary this
ratio, asking what follows if maintenance is primary.

\medskip

\noindent\textbf{Ecology as substrate vs.\ ecology as externality.}
Contemporary economic organization treats ecological processes as
externalities — inputs and outputs that are real but not directly
represented in the accounting. The proposals vary this by imagining
systems in which ecological processes are the primary organizing
principle.

\medskip

\noindent\textbf{Persistence vs.\ production.} Contemporary value
systems reward production — the creation of new objects, new
information, new services — and systematically undervalue persistence:
the maintenance of existing knowledge, social bonds, biological systems,
and physical infrastructure. The proposals vary this by imagining systems
organized around persistence as the primary value.

\medskip

\noindent\textbf{Accessibility vs.\ location.} Contemporary geography
treats location as a primary determinant of social possibility. The
proposals vary this by imagining systems in which what matters is
accessibility — the range of futures reachable from a given state —
rather than fixed spatial coordinates.

\medskip

\noindent\textbf{Coordination as first-class design vs.\ coordination as
emergent.} Contemporary technical systems typically treat coordination
as an emergent property of individual optimization. The proposals vary
this by imagining systems in which coordination mechanisms are
explicitly designed rather than assumed to arise spontaneously.

\medskip

These five parameters — maintenance, ecology, persistence,
accessibility, coordination — define the invariant structure that the
individual proposals orbit. The armillary sphere is not any one ring.
The armillary sphere is the five-dimensional coordinate system within
which civilization's design choices can be situated and varied.

These parameters can be treated as the objects of formal study.

\begin{definition}[Civilizational Invariant]
A quantity $I: \mathcal{C} \to \mathbb{R}$ is a \emph{civilizational
invariant} under a family of transformations $\{T_i\}$ if
\[
  I(T_i(C)) = I(C) \quad \text{for all } i.
\]
An invariant is a structural feature of civilizational configurations
that persists across the variation introduced by armillary proposals.
\end{definition}

\begin{proposition}[Invariant Discovery Principle]
The purpose of an armillary proposal is not to identify an optimal
configuration $C^\star$ but to discover quantities $I$ that remain
invariant under broad families of transformations. What is preserved
under variation is more informative about the deep structure of
civilization than what is changed.
\end{proposition}

The reachability volume connects this to the admissibility analysis
developed elsewhere. For a configuration $C$, let $R(C)$ denote the
set of future configurations reachable under admissible transitions,
and define:

\begin{definition}[Reachability Volume]
\[
  \Omega(C) = \mathrm{Vol}(R(C)).
\]
\end{definition}

\begin{proposition}[Reachability Expansion]
A successful armillary proposal increases understanding when it reveals
directions $v$ in $T_C\mathcal{C}$ such that
\[
  \frac{d}{d\epsilon}\,\Omega(C + \epsilon v) > 0.
\]
That is: moving in direction $v$ expands the set of reachable futures,
not merely the set of imagined ones.
\end{proposition}

The Armillary Principle can now be stated precisely:

\begin{proposition}[Armillary Principle]
Let $T$ be a coordinate transformation proposal. Its value is
determined not by $\|T(C) - C\|$ — the magnitude of the imagined
departure from present reality — but by $\Delta I$ and $\Delta\Omega$:
the new invariants it reveals and the reachable possibilities it
exposes.
\end{proposition}

The individual proposals are not answers to the question of how to
organize civilization. They are instruments for making the question
more precise.

% ---------------------------------------------------------------
\section{What the Satire Reveals}

The satirical reading of speculative proposals is not simply an error.
It is revealing. The indecision that characterizes the best satirical
commentary on this kind of work — the oscillation between mockery and
fascination, the inability to decide whether one is looking at a plan, a
joke, or something else — is a symptom of encountering an object that
does not fit available categories.

When a reader cannot decide whether a proposal is serious or absurd,
they are often experiencing a genuine category problem. The proposal is
occupying the intermediate space between blueprint and metaphor that the
armillary sphere represents, and the reader's interpretive framework
offers only two slots: literal or figurative. Neither slot fits.

The most productive response to this discomfort is not to force the
proposal into one of the available categories but to recognize the
discomfort as diagnostic. Something is being exposed. The familiar
coordinate system is insufficient for locating this object. A new
coordinate is needed.

That recognition — the experience of a familiar framework proving
insufficient — is precisely what the thought experiment and the
coordinate transformation proposal are designed to provoke. The
discomfort is not a failure of the proposal. It is its purpose.

When someone describes forty billion floating habitat modules and the
description sounds absurd, the absurdity is a signal. It is revealing
that the parameters being varied — the substrate of settlement, the
fixity of location, the centralization of infrastructure — are so
deeply embedded that varying them sounds like nonsense rather than like
a question. The satire works because the assumptions are invisible.
The proposal works because it makes them visible.

The satirist and the speculative designer are, in this moment, pointing
at the same thing from opposite directions.

% ---------------------------------------------------------------
\section*{Conclusion}

An armillary sphere does not tell you where to sail. It tells you how
to think about where you are. Its rings are not real celestial circles;
they are instruments for navigating real celestial relationships. The
sphere is not a model of the world; it is a model of the coordinate
system within which the world can be understood and traversed.

Speculative proposals that vary the fundamental parameters of
civilization — that treat maintenance as primary, ecology as substrate,
persistence as value, accessibility as geography, coordination as
design — are armillary spheres. They are not blueprints for a future
civilization. They are instruments for thinking about what civilization
is, what it assumes, and which of those assumptions are contingent
rather than necessary.

The question these proposals collectively probe is not: what should we
build? It is: what becomes visible when we release the parameters that
habituation has rendered invisible?

Rotating civilization around the axis of maintenance: what was hidden
becomes visible about the costs of disposability.

Rotating civilization around the axis of ecology: what was hidden becomes
visible about the subsidy of extraction.

Rotating civilization around the axis of persistence: what was hidden
becomes visible about the systematic undervaluation of continuity.

Rotating civilization around the axis of accessibility: what was hidden
becomes visible about the political geography of possibility.

None of these rotations produce a plan. All of them produce a question.
The questions are the point. The proposals are the instrument for
generating the questions. The armillary sphere is not the Earth.

It helps you find your position in relation to it.

% ---------------------------------------------------------------
\begin{thebibliography}{99}

% Thought experiments in philosophy and physics
\bibitem{sorensen1992}
Sorensen, R. A. (1992).
\textit{Thought Experiments}.
Oxford University Press.

\bibitem{brown1991}
Brown, J. R. (1991).
\textit{The Laboratory of the Mind: Thought Experiments in the Natural Sciences}.
Routledge.

\bibitem{mach1883}
Mach, E. (1883).
\textit{The Science of Mechanics}.
[English trans.: Open Court, 1960.]

\bibitem{norton1996}
Norton, J. D. (1996).
Are Thought Experiments Just What You Thought?
\textit{Canadian Journal of Philosophy}, 26(3), 333--366.

% Coordinate systems and the history of navigation
\bibitem{neugebauer1975}
Neugebauer, O. (1975).
\textit{A History of Ancient Mathematical Astronomy}.
Springer.

\bibitem{dreyer1906}
Dreyer, J. L. E. (1906).
\textit{History of the Planetary Systems from Thales to Kepler}.
Cambridge University Press.

% Design theory and provocation
\bibitem{dunne2013}
Dunne, A., \& Raby, F. (2013).
\textit{Speculative Everything: Design, Fiction, and Social Dreaming}.
MIT Press.

\bibitem{schon1983}
Sch\"{o}n, D. A. (1983).
\textit{The Reflective Practitioner: How Professionals Think in Action}.
Basic Books.

% Philosophy of science and models
\bibitem{cartwright1983}
Cartwright, N. (1983).
\textit{How the Laws of Physics Lie}.
Oxford University Press.

\bibitem{giere1988}
Giere, R. N. (1988).
\textit{Explaining Science: A Cognitive Approach}.
University of Chicago Press.

\bibitem{morgan2012}
Morgan, M. S. (2012).
\textit{The World in the Model: How Economists Work and Think}.
Cambridge University Press.

% Maps, representations, and territory
\bibitem{korzybski1933}
Korzybski, A. (1933).
\textit{Science and Sanity}.
International Non-Aristotelian Library.

\bibitem{borges1946}
Borges, J. L. (1946).
Del rigor en la ciencia.
\textit{Los Anales de Buenos Aires}, 1(3).
[English: On Exactitude in Science.]

\bibitem{bateson1972}
Bateson, G. (1972).
\textit{Steps to an Ecology of Mind}.
Chandler.

% Utopia, design failure, and social systems
\bibitem{jacobs1961}
Jacobs, J. (1961).
\textit{The Death and Life of Great American Cities}.
Random House.

\bibitem{scott1998}
Scott, J. C. (1998).
\textit{Seeing Like a State: How Certain Schemes to Improve the Human
Condition Have Failed}.
Yale University Press.

\bibitem{winner1980}
Winner, L. (1980).
Do Artifacts Have Politics?
\textit{Daedalus}, 109(1), 121--136.

% Reachability and admissibility
\bibitem{ostrom1990}
Ostrom, E. (1990).
\textit{Governing the Commons: The Evolution of Institutions for
Collective Action}.
Cambridge University Press.

\bibitem{meadows2008}
Meadows, D. H. (2008).
\textit{Thinking in Systems: A Primer}.
Chelsea Green.

% Maintenance and repair
\bibitem{graham2001}
Graham, S., \& Thrift, N. (2007).
Out of Order: Understanding Repair and Maintenance.
\textit{Theory, Culture \& Society}, 24(3), 1--25.

\bibitem{jackson2014}
Jackson, S. J. (2014).
Rethinking Repair.
In T. Gillespie, P. Boczkowski, \& K. Foot (Eds.),
\textit{Media Technologies: Essays on Communication, Materiality, and Society}.
MIT Press.

% Ecology and design
\bibitem{mcdounough2002}
McDonough, W., \& Braungart, M. (2002).
\textit{Cradle to Cradle: Remaking the Way We Make Things}.
North Point Press.

\bibitem{odum1971}
Odum, H. T. (1971).
\textit{Environment, Power, and Society}.
Wiley-Interscience.

% Complexity, emergence, and civilization
\bibitem{holland1995}
Holland, J. H. (1995).
\textit{Hidden Order: How Adaptation Builds Complexity}.
Addison-Wesley.

\bibitem{simon1962}
Simon, H. A. (1962).
The Architecture of Complexity.
\textit{Proceedings of the American Philosophical Society}, 106(6), 467--482.

\end{thebibliography}

% ---------------------------------------------------------------
\end{document}
